% bismuto: schema delle bande vicino a E_F (Hofmann fig. 1b della nota sul bismuto)
\begin{tikzpicture}
\begin{axis}[width=8.6cm, height=5.4cm, xmin=-1, xmax=3, ymin=-1.3, ymax=0.9, axis lines=box, xtick={0,2}, xticklabels={$T$,$L$},
  ytick={0}, yticklabels={$\EF$}, ylabel={$E$}, samples=101]
  \addplot[very thick, color=vB, domain=-0.75:0.75] {0.2-2*x^2};
  \addplot[very thick, color=vA, domain=1.35:2.65] {-0.2+1.2*abs(x-2)};
  \addplot[very thick, color=vB, domain=1.35:2.65] {-0.55-1.2*abs(x-2)};
  \addplot[black!50, dashed] coordinates {(-1,0) (3,0)};
  \addplot[color=hole!70, fill=hole!35] coordinates {(-0.316,0) (-0.2,0.12) (0,0.2) (0.2,0.12) (0.316,0)} -- cycle;
  \addplot[color=exc!70, fill=exc!35] coordinates {(1.833,0) (2,-0.2) (2.167,0)} -- cycle;
  \draw[<->, >=Stealth] (axis cs:1.0,-0.2) -- (axis cs:1.0,0.2) node[above, font=\scriptsize] {$\approx\SI{40}{meV}$};
  \draw[black!40, densely dotted] (axis cs:0,0.2) -- (axis cs:1.0,0.2); \draw[black!40, densely dotted] (axis cs:1.0,-0.2) -- (axis cs:2,-0.2);
  \node[font=\scriptsize, color=hole] at (axis cs:0,0.45) {lacune};
  \node[font=\scriptsize, color=exc!70!black] at (axis cs:2.0,-0.36) {elettroni};
  \node[font=\scriptsize, color=vB] at (axis cs:-0.2,-1.05) {«valenza»};
  \node[font=\scriptsize, color=vA] at (axis cs:2.4,0.75) {«conduzione»};
\end{axis}
\end{tikzpicture}
